\documentclass{article}
\input{../../../lib.sty}

\title{\texttt{fn log\_mean\_hi}}
\author{Michael Shoemate}

\begin{document}
\maketitle

Proves soundness of \texttt{fn log\_mean\_hi} in \texttt{mod.rs}.

\section{Hoare Triple}
\subsection*{Precondition}
\texttt{mean} $> 0$ and \texttt{mean} is finite.

\subsection*{Pseudocode}
\lstinputlisting[language=Python,firstline=2,escapechar=|]{./pseudocode/log_mean_hi.py}

\subsection*{Postcondition}
\begin{theorem}
\label{postcondition}
For any setting of the input parameter \texttt{mean} such that the given precondition holds,
\texttt{log\_mean\_hi} either returns \texttt{Err(e)},
or returns \texttt{Ok(out)} where
\[
\texttt{out} \ge \log(\texttt{mean}).
\]
\end{theorem}

\section{Proof}
Assume the precondition is met.

\begin{lemma}
\label{err-e}
\texttt{log\_mean\_hi} only returns \texttt{Err(e)} if the guard on line \ref{line:guard} fails
or if a later fallible helper call returns \texttt{Err(e)}.
\end{lemma}

\begin{proof}
The function explicitly returns an error when the guard on line \ref{line:guard} fails.
Under the precondition, this branch is not taken.
The only remaining fallible operations are the calls on lines \ref{line:minus_one_hi}
and \ref{line:ln1p_hi}.
\end{proof}

\begin{lemma}
\label{ok-out}
If the outcome of \texttt{log\_mean\_hi} is \texttt{Ok(out)}, then
\[
\texttt{out} \ge \log(\texttt{mean}).
\]
\end{lemma}

\begin{proof}
By the proof definition of \texttt{inf\_sub},
\[
\texttt{mean\_minus\_one\_hi} \ge \texttt{mean} - 1.
\]
Since \texttt{mean} $> 0$, both arguments to $\ln(1+\cdot)$ lie in $(-1,\infty)$.
By the proof definition of \texttt{inf\_ln\_1p},
\[
\texttt{out}
=
\texttt{mean\_minus\_one\_hi.inf\_ln\_1p()}
\ge
\ln((\texttt{mean} - 1) + 1)
= \log(\texttt{mean}).
\]
\end{proof}

\begin{proof}[Proof of Theorem~\ref{postcondition}]
Holds by Lemma~\ref{err-e} and Lemma~\ref{ok-out}.
\end{proof}

\end{document}
